\section{模型评价与推广}

\subsection{模型的优点}

\textbf{（1）预测模型由数据规律驱动，可解释性强。} 本文并未直接套用通用的时间序列预测方法，
而是先对附件 2 的全年数据作统计分析，发现小区负载具有显著的星期型双状态结构
（周日至周四日均负载约 124.0\,MWh，周五、周六仅约 78.3\,MWh，比值 0.63），
据此选择“上周同星期实际值”作为日前预测的基础，从而避免了持续性预测在星期状态切换日
（如 12 月 20 日周六至 21 日周日）预测失效的问题；光伏出力则具有强季节性与较高的日间持续性
（日间自相关系数 0.964），支持使用近期历史数据。模型变量与参数均有明确的实际含义。

\textbf{（2）关键参数具有理论依据而非经验试凑。} 5 倍电价惩罚使购电决策具有强不对称性，
其结构与报童模型一致，临界分位数 $(5-1)/5=0.8$ 要求计划购电覆盖净负荷约 80\% 分位点，
据此设置的安全余量在全年扫描中形成平坦平台（$\gamma$ 在 $[1.03,1.07]$ 内总费用变化不足 1\%），
且分段验证结论一致，说明参数取值稳健、模型未发生过拟合。

\textbf{（3）模型分层递进，四个问题共用同一优化内核。} 从问题一的日循环线性规划，
到问题二“预测—安全余量—日前计划—日内纠偏—跨日滚动”，
再到问题三“预报标定—残差修正—双向调整—末端备用”以及问题四的价格项替换，
模型结构始终一致、各层功能清晰，便于横向比较不同策略的经济性。

\textbf{（4）求解可靠、结果可复现且具有经济含义。} 日前计划与调整计划均为线性规划，
可求得全局最优；日内纠偏规则透明、物理可行。全链路数值校验（功率平衡、储能越限、跨日连续、
费用分解、调整 LP 复现）在 334 天滚动仿真中零违规，且数值求解输出、结果文件与模型说明三方一致。
在效率上，单次全年仿真（334 次日前规划 $+$ 约 1000 次调整规划）可在数分钟内完成，
使余量参数、残差权重与末端备用的大规模扫描寻优成为可能。
从结果看，储能设备的年价值为 439.2 万元（取消储能后费用上升 23.2\%）；
问题三把 5 倍紧急购电转化为 0.5 倍与 1.5 倍的调整结算后，全年费用较问题二下降 5.9\%，
紧急购电量由 14.51 万\,kWh 降至 5.59 万\,kWh，与费率结构改善的机理完全吻合。

\subsection{模型的不足}

\textbf{（1）预测仍为点预测加比例余量。} 本文未建立负载与光伏的联合概率分布或场景集，
余量系数按全天统一设置、未按时段细化，因此余量策略存在固有代价：
问题二全年弃电约 326 万\,kWh，其中“已付费却被迫丢弃”的电量下限约 81 万\,kWh。
若采用分布鲁棒优化\cite{bertsimas2013}或分时段余量，这部分损失有望进一步压缩。

\textbf{（2）日内调度采用贪心规则。} 储能纠偏按固定优先级执行，未在日内逐时段重新优化，
与模型预测控制（MPC）\cite{mayne2000}相比仍留有一定优化空间（估计在 1\%$\sim$2\% 量级）。

\textbf{（3）电价不确定性与参数寻优方式尚不完善。} 问题四假设计划制定时当日逐时电价已知，
而实际电力市场中电价本身也需预测，若电价未知则应引入电价场景或不确定集，
将模型扩展为两阶段随机规划，本文仅将其作为推广方向；
此外，余量系数、残差更新权重与末端备用水平均以全年数据寻优确定，
虽已通过平台分析与分段验证排除过拟合，但严格意义上仍应采用滚动交叉验证。

\textbf{（4）目标函数未考虑储能的循环损耗。} 模型以购电费用最小为唯一目标，
未计入充放电次数带来的寿命损耗。若引入单位循环成本，储能周转量会相应下降，
总费用与最优策略可能发生变化。

\subsection{模型的推广}

本文模型的核心思想——\textbf{在不对称惩罚下用保守余量制定日前计划、用日内滚动调整校正预测偏差}——
并不依赖具体的微网场景，可推广到其他“日前承诺、实时考核”类型的决策问题。

\textbf{（1）同类电力系统与设备。} 风电场与光伏电站的日前出力申报、独立储能电站的峰谷套利、
园区综合能源系统的多能互补调度、需求响应中可中断负荷的报价等问题，
均具有“计划量按合同结算、偏差按惩罚电价结算”的结构，可直接沿用本文的两阶段框架。
推广时只需将本文的惩罚倍数（5 倍、1.5 倍、0.5 倍）替换为对应场景的考核系数，
将储能模型替换为电池、储热或储氢设备的容量—功率—效率参数即可。

\textbf{（2）数据与机制变化下的调整。} 若预报发布更频繁（如每 15\,min 一次），
可相应加密调整时刻的数量；若电价未知，可将固定电价替换为价格场景集或不确定集，
把模型扩展为两阶段随机规划或鲁棒优化；若充放电效率随荷电状态变化，
则可将当前的线性储能模型替换为分段线性模型。

\textbf{（3）方法层面的改进方向。} 预测环节可用 LSTM、Transformer 等机器学习模型
替代本文的 MOS 线性回归，或引入概率预测方法\cite{hong2016}，并与光伏出力的物理机理模型结合形成混合预测；
日内调度可由贪心规则升级为模型预测控制，使储能动作在滚动时域内持续优化；
此外还可引入储能寿命模型，把循环损耗纳入目标函数，从而得到更贴近实际运行的调度策略。
